\documentclass[a4paper,11pt]{article}
\usepackage[utf8]{inputenc}
\usepackage[T1]{fontenc}
\usepackage[ngerman]{babel}
\usepackage{lmodern}
\usepackage{microtype}
\usepackage{amsmath,amssymb,bm,mathtools}
\usepackage{geometry}
\geometry{a4paper, left=2.4cm, right=2.4cm, top=2.0cm, bottom=2.2cm}
\usepackage{booktabs,multirow,array}
\usepackage{xcolor}
\definecolor{bokug}{RGB}{0,135,60}
\usepackage{tikz}
\usetikzlibrary{patterns,arrows.meta,calc,decorations.pathmorphing,
                 decorations.markings}
\usepackage{fancyhdr}
\pagestyle{fancy}
\renewcommand{\headrulewidth}{0.4pt}
\renewcommand{\footrulewidth}{0.4pt}
\usepackage{enumitem}
\setlength{\parindent}{0pt}
\setlength{\parskip}{4pt}
\begin{document}

\fancyhead[L]{\small Musterlösung\quad VU Mechanik -- LAWI 100515 (PI)\quad SS 2026}
\fancyhead[R]{\small 2.~Übungsklausur (26.06.2026)\quad Gruppe~A}
\fancyfoot[C]{\small Musterlösung}
\fancyfoot[R]{\small Seite~\thepage}

\begin{center}
  {\Large\bfseries Musterlösung}\\[4pt]
  {\large VU Mechanik -- LAWI 100515 (PI)\quad SS 2026}\\[2pt]
  {\normalsize 2.~Übungsklausur (26.06.2026) -- Gruppe~A}\\[2pt]
  {\normalsize \textbf{Systemwerte:} $L = 2$\,m,\quad $q = 3$\,kN/m,\quad $P = 20$\,kN}
\end{center}

\section*{1.~Beispiel \hfill (70\,\%)}


\begin{center}
\begin{tikzpicture}[scale=3.2, >=stealth, line cap=round, line join=round,
    thick, font=\small]

  %--- Node coordinates (normalized: 1 unit = L) ----------------------------
  \coordinate (A)   at (0.00, 0.00);
  \coordinate (C)   at (1.00, 0.00);
  \coordinate (G)   at (2.00, 0.00);
  \coordinate (D)   at (1.00,-0.50);
  \coordinate (E)   at (1.50,-0.50);
  \coordinate (Fnd) at (1.50,-1.00);
  \coordinate (B)   at (2.00,-1.00);

  %--- Fixed support (Einspannung) at A ------------------------------------
  \fill[pattern=north east lines, pattern color=black!55]
        (-0.14,-0.22) rectangle (0.00, 0.22);
  \draw[very thick] (0,-0.22) -- (0, 0.22);

  %--- Beam members [1] and [2] --------------------------------------------
  \draw[very thick] (A) -- (G);

  %--- Truss members -------------------------------------------------------
  \draw (C) -- (D);         % [3]
  \draw (C) -- (E);         % [4]
  \draw (D) -- (E);         % [5]
  \draw (D) -- (Fnd);       % [6]
  \draw (E) -- (Fnd);       % [7]
  \draw (E) -- (B);         % [8]
  \draw (Fnd) -- (B);       % [9]

  %--- Hinge at C (open circle) -------------------------------------------
  \filldraw[fill=white, draw=black, thin] (C) circle (1.3pt);

  %--- Roller support at B (verschiebliches Lager, vertical reaction) ------
  \draw (B) -- ++(0.09,-0.13) -- ++(-0.18, 0) -- cycle;
  \foreach \dx in {-0.09, -0.03, 0.03, 0.09}
    \filldraw[fill=white, draw=black, thin]
              ($(B)+(\dx,-0.13)$) circle (1.1pt);
  \draw[thick] ($(B)+(-0.13,-0.16)$) -- ++(0.26, 0);
  \fill[pattern=north east lines, pattern color=black!40]
       ($(B)+(-0.13,-0.19)$) rectangle ++(0.26,-0.05);

  %--- Distributed load q (downward arrows above beam) --------------------
  \foreach \xx in {0.10,0.30,0.50,0.70,0.90,1.10,1.30,1.50,1.70,1.90}
    \draw[->, thin, blue!75!black] (\xx, 0.32) -- (\xx, 0.04);
  \draw[thin, blue!75!black] (0.05, 0.32) -- (1.95, 0.32);
  \node[blue!75!black, right, inner sep=1pt] at (1.97, 0.32) {$q$};

  %--- Point load P at D (horizontal, rightward) --------------------------
  \draw[->, very thick, red!70!black]
        ($(D)+(-0.48,0)$) -- ($(D)+(-0.06,0)$);
  \node[red!70!black, left, inner sep=1pt] at ($(D)+(-0.50,0)$) {$P$};

  %--- Node labels ---------------------------------------------------------
  \node[left,  xshift=-2pt] at (A)   {\textbf{A}};
  \node[above, yshift= 3pt] at (C)   {\textbf{C}};
  \node[above right]        at (G)   {\textbf{G}};
  \node[left,  xshift=-2pt] at (D)   {\textbf{D}};
  \node[above right]        at (E)   {\textbf{E}};
  \node[right, xshift= 2pt] at (Fnd) {\textbf{F}};
  \node[right, xshift= 2pt] at (B)   {\textbf{B}};

  %--- Member labels -------------------------------------------------------
  \node[above]        at (0.50, 0.01) {[1]};
  \node[above]        at (1.50, 0.01) {[2]};
  \node[left]         at (1.00,-0.25) {[3]};
  \node[above right, xshift=-1pt] at (1.23,-0.23) {[4]};
  \node[above]        at (1.25,-0.50) {[5]};
  \node[below left, xshift=1pt]   at (1.23,-0.77) {[6]};
  \node[right]        at (1.50,-0.75) {[7]};
  \node[above right, xshift=-1pt] at (1.73,-0.72) {[8]};
  \node[below]        at (1.75,-1.00) {[9]};

  %--- Dimension lines -----------------------------------------------------
  \draw[<->, thin, gray] (0.00,-1.35) -- (1.00,-1.35)
        node[midway, below, font=\footnotesize] {$L$};
  \draw[<->, thin, gray] (1.00,-1.35) -- (1.50,-1.35)
        node[midway, below, font=\footnotesize] {$L/2$};
  \draw[<->, thin, gray] (1.50,-1.35) -- (2.00,-1.35)
        node[midway, below, font=\footnotesize] {$L/2$};
  \draw[<->, thin, gray] (2.30, 0.00) -- (2.30,-0.50)
        node[midway, right, font=\footnotesize] {$L/2$};
  \draw[<->, thin, gray] (2.30,-0.50) -- (2.30,-1.00)
        node[midway, right, font=\footnotesize] {$L/2$};

\end{tikzpicture}
\end{center}


\subsection*{1.1\quad Statische Bestimmtheit}

Das Tragwerk besteht aus \textbf{zwei starren Scheiben}: dem Biegeträger A--C--G
(Stäbe~[1],\,[2]) und dem Fachwerk (Knoten C,D,E,F,B; Stäbe~[3]--[9]).

\textit{Gesamtsystem:} 2 Scheiben $\Rightarrow$ $2 \times 3 = 6$ Gleichgewichtsbedingungen.
Unbekannte: Auflager A (Einspannung) $= 3$, Auflager B (verschiebliches Lager) $= 1$,
Gelenkverbindung C zwischen Träger und Fachwerk $= 2$. Summe $= 6$.
\[
  n = \underbrace{3}_{A} + \underbrace{1}_{B} + \underbrace{2}_{C} - 2 \cdot 3
    = 6 - 6 = 0
  \quad\Rightarrow\quad \boxed{n = 0: \text{statisch bestimmt}}
\]

\subsection*{1.2\quad Auflagerreaktionen (allgemein als Funktion von $q$, $P$, $L$)}

\textbf{Fachwerk als Teilsystem}
(Gelenkskräfte $C_x$, $C_y$ vom Träger auf das Fachwerk; $B_V$ in B; $P$ in D):
\begin{align*}
  \sum M_C = 0:\quad & P \cdot \tfrac{L}{2} + B_V \cdot L = 0
    &&\Rightarrow\quad \boxed{B_V = -\tfrac{P}{2}} \\
  \sum F_x = 0:\quad & C_x + P = 0
    &&\Rightarrow\quad C_x = -P \\
  \sum F_y = 0:\quad & C_y + B_V = 0
    &&\Rightarrow\quad C_y = +\tfrac{P}{2}
\end{align*}

\textbf{Gesamtsystem} (Gleichlast $Q = q \cdot 2L$ mit Angriffspunkt $x = L$):
\begin{align*}
  \sum F_x = 0:\quad & A_H + P = 0
    &&\Rightarrow\quad \boxed{A_H = -P} \\
  \sum F_y = 0:\quad & A_V + B_V - 2qL = 0
    &&\Rightarrow\quad \boxed{A_V = 2qL + \tfrac{P}{2}} \\
  \sum M_A = 0:\quad & M_A - 2qL \cdot L + P\cdot\tfrac{L}{2} + B_V \cdot 2L = 0
    &&\Rightarrow\quad \boxed{M_A = 2qL^2 + \tfrac{PL}{2}}
\end{align*}

\subsection*{1.3\quad Zahlenwerte}

\begin{center}
\begin{tabular}{cccc}
\toprule
$A_V$ & $A_H$ & $M_A$ & $B_V$ \\
\midrule
$2qL + \tfrac{P}{2} = \mathbf{22.00}$\,kN
& $-P = \mathbf{-20.00}$\,kN
& $2qL^2 + \tfrac{PL}{2} = \mathbf{44.00}$\,kNm
& $-\tfrac{P}{2} = \mathbf{-10.00}$\,kN \\
\bottomrule
\end{tabular}
\end{center}

Vorzeichen: positiv $\uparrow$ für $A_V$, $B_V$; positiv $\rightarrow$ für $A_H$;
CCW für $M_A$. Negativ $\Rightarrow$ tatsächliche Richtung entgegengesetzt.

\textit{Kontrolle} $\sum M_B = 0$:
$M_A - A_H \cdot L - 2A_V \cdot L + 2qL^2 - \tfrac{PL}{2} = 0$\;
\checkmark

\subsection*{1.4\quad Schnittgrößen in den Stäben [1] und [2]}

Am Knoten C wirkt auf den Träger die Fachwerkkraft
$(-C_x,\,-C_y) = (+P,\,-\tfrac{P}{2})$:
$P$ nach rechts und $\tfrac{P}{2}$ nach unten.
Laufvariable $x$ jeweils ab Stabanfang.

\textbf{Stab~[1] (A\,$\to$\,C), $0 \le x \le L$:}
\begin{align*}
  N_1(x) &= P = \text{const} \\
  V_1(x) &= \Bigl(2qL + \tfrac{P}{2}\Bigr) - qx \\
  M_1(x) &= \Bigl(2qL + \tfrac{P}{2}\Bigr)x - \tfrac{q}{2}x^2
            - \Bigl(2qL^2 + \tfrac{PL}{2}\Bigr)
\end{align*}

\textbf{Stab~[2] (C\,$\to$\,G), $0 \le x \le L$:}
\begin{align*}
  N_2(x) &= 0 \\
  V_2(x) &= q(L - x) \\
  M_2(x) &= -\tfrac{q}{2}(L - x)^2
\end{align*}

Am Übergang C: Querkraftsprung $-\tfrac{P}{2}$ (vertikale Fachwerkkraft);
Normalkraftsprung $-P$ (horizontale Fachwerkkraft); Moment stetig, jedoch mit Knick.

\medskip
\begin{center}
\begin{tabular}{l cccccc}
\toprule
 & \multicolumn{3}{c}{\textbf{Stabanfang}} & \multicolumn{3}{c}{\textbf{Stabende}} \\
\cmidrule(lr){2-4}\cmidrule(lr){5-7}
 & $M$\,[kNm] & $V$\,[kN] & $N$\,[kN]
 & $M$\,[kNm] & $V$\,[kN] & $N$\,[kN] \\
\midrule
Stab~[1] (A$\to$C) & $-44.00$ & $22.00$ & $20.00$ & $-6.00$ & $16.00$ & $20.00$ \\[3pt]
Stab~[2] (C$\to$G) & $-6.00$ & $6.00$ & $0.00$ & $-0.00$ & $0.00$ & $0.00$ \\
\bottomrule
\end{tabular}
\end{center}

\subsection*{1.5\quad Schnittgrößenverläufe}


%─────────────────────────────────────────────────────────────────────────────
\subsubsection*{Schnittgrößenverläufe (auf der Zugseite; neg.\,M nach oben)}

\begin{center}
\begin{tikzpicture}[xscale=3.2, yscale=3.2, >=stealth, thick, font=\small]

  %--- Axis ticks and node labels for ALL diagrams -------------------------
  % We stack 3 diagrams vertically: M at y-offset=0, V at y-offset=-2, N at y-offset=-4

  %====================================================================
  % MOMENT  M(x)
  %====================================================================
  \begin{scope}[yshift= 0cm]
    % Baseline
    \draw[->] (-0.1, 0) -- (2.15, 0) node[right] {$x$};
    \node[left] at (0,0)   {A};
    \node[above] at (1,0)  {C};
    \node[right] at (2,0)  {G};
    \draw[thin, dashed] (1,0) -- (1, 0.1023);

    % M in Stab [1]: filled area (tension side = above)
    \filldraw[fill=blue!15, draw=blue!60!black, thick]
      (0,0) -- plot coordinates { (0.00000,0.75000) (0.03333,0.72511) (0.06667,0.70045) (0.10000,0.67602) (0.13333,0.65182) (0.16667,0.62784) (0.20000,0.60409) (0.23333,0.58057) (0.26667,0.55727) (0.30000,0.53420) (0.33333,0.51136) (0.36667,0.48875) (0.40000,0.46636) (0.43333,0.44420) (0.46667,0.42227) (0.50000,0.40057) (0.53333,0.37909) (0.56667,0.35784) (0.60000,0.33682) (0.63333,0.31602) (0.66667,0.29545) (0.70000,0.27511) (0.73333,0.25500) (0.76667,0.23511) (0.80000,0.21545) (0.83333,0.19602) (0.86667,0.17682) (0.90000,0.15784) (0.93333,0.13909) (0.96667,0.12057) (1.00000,0.10227) } -- (1,0) -- cycle;

    % M in Stab [2]: filled area
    \filldraw[fill=blue!15, draw=blue!60!black, thick]
      (1,0) -- plot coordinates { (1.00000,0.10227) (1.03333,0.09557) (1.06667,0.08909) (1.10000,0.08284) (1.13333,0.07682) (1.16667,0.07102) (1.20000,0.06545) (1.23333,0.06011) (1.26667,0.05500) (1.30000,0.05011) (1.33333,0.04545) (1.36667,0.04102) (1.40000,0.03682) (1.43333,0.03284) (1.46667,0.02909) (1.50000,0.02557) (1.53333,0.02227) (1.56667,0.01920) (1.60000,0.01636) (1.63333,0.01375) (1.66667,0.01136) (1.70000,0.00920) (1.73333,0.00727) (1.76667,0.00557) (1.80000,0.00409) (1.83333,0.00284) (1.86667,0.00182) (1.90000,0.00102) (1.93333,0.00045) (1.96667,0.00011) (2.00000,0.00000) } -- (2,0) -- cycle;

    % Value labels
    \node[right, blue!70!black] at (0, 0.7500)
          {$M=-44.0$\,kNm};
    \node[right, blue!70!black] at (1, 0.1023)
          {$M=-6.0$\,kNm};
    \node[above] at (1,-0.05) {\footnotesize C};

    % Axis label
    \node[rotate=90, blue!70!black] at (-0.35, 0.3750) {$M(x)$};
  \end{scope}

  %====================================================================
  % SHEAR FORCE  V(x)
  %====================================================================
  \begin{scope}[yshift=-2.8cm]
    \draw[->] (-0.1, 0) -- (2.15, 0) node[right] {$x$};
    \node[left]  at (0,0)  {A};
    \node[above] at (1,-0.02) {C};
    \node[right] at (2,0)  {G};

    % V in Stab [1]
    \filldraw[fill=red!15, draw=red!60!black, thick]
      (0,0) -- plot coordinates { (0.00000,0.70000) (0.03333,0.69364) (0.06667,0.68727) (0.10000,0.68091) (0.13333,0.67455) (0.16667,0.66818) (0.20000,0.66182) (0.23333,0.65545) (0.26667,0.64909) (0.30000,0.64273) (0.33333,0.63636) (0.36667,0.63000) (0.40000,0.62364) (0.43333,0.61727) (0.46667,0.61091) (0.50000,0.60455) (0.53333,0.59818) (0.56667,0.59182) (0.60000,0.58545) (0.63333,0.57909) (0.66667,0.57273) (0.70000,0.56636) (0.73333,0.56000) (0.76667,0.55364) (0.80000,0.54727) (0.83333,0.54091) (0.86667,0.53455) (0.90000,0.52818) (0.93333,0.52182) (0.96667,0.51545) (1.00000,0.50909) } -- (1,0) -- cycle;

    % V in Stab [2]
    \filldraw[fill=red!15, draw=red!60!black, thick]
      (1,0) -- plot coordinates { (1.00000,0.19091) (1.03333,0.18455) (1.06667,0.17818) (1.10000,0.17182) (1.13333,0.16545) (1.16667,0.15909) (1.20000,0.15273) (1.23333,0.14636) (1.26667,0.14000) (1.30000,0.13364) (1.33333,0.12727) (1.36667,0.12091) (1.40000,0.11455) (1.43333,0.10818) (1.46667,0.10182) (1.50000,0.09545) (1.53333,0.08909) (1.56667,0.08273) (1.60000,0.07636) (1.63333,0.07000) (1.66667,0.06364) (1.70000,0.05727) (1.73333,0.05091) (1.76667,0.04455) (1.80000,0.03818) (1.83333,0.03182) (1.86667,0.02545) (1.90000,0.01909) (1.93333,0.01273) (1.96667,0.00636) (2.00000,0.00000) } -- (2,0) -- cycle;

    % Jump at C (vertical dashed line)
    \draw[dashed, thin] (1, 0.5091) -- (1, 0.1909);

    % Labels
    \node[right, red!70!black] at (0,  0.7000) {$V=22.0$\,kN};
    \node[right, red!70!black] at (1,  0.5091) {$V=16.0$\,kN};
    \node[left,  red!70!black] at (1,  0.1909) {$V=6.0$\,kN};
    \node[rotate=90, red!70!black] at (-0.35, 0.3500) {$V(x)$};
  \end{scope}

  %====================================================================
  % AXIAL FORCE  N(x)
  %====================================================================
  \begin{scope}[yshift=-5.6cm]
    \draw[->] (-0.1, 0) -- (2.15, 0) node[right] {$x$};
    \node[left]  at (0,0)  {A};
    \node[above] at (1,-0.02) {C};
    \node[right] at (2,0)  {G};

    % N in Stab [1]: constant positive (tension, draw below baseline)
    \filldraw[fill=green!15, draw=green!50!black, thick]
      (0,0) -- (0,-0.6500) -- (1,-0.6500) -- (1,0) -- cycle;

    % N in Stab [2]: zero
    % (nothing to draw)
    \draw[green!50!black, thick] (1,0) -- (2,0);

    % Labels
    \node[right, green!40!black] at (0,  -0.6500)
          {$N=+20.0$\,kN\,(Zug)};
    \node[rotate=90, green!40!black] at (-0.35, -0.3250) {$N(x)$};
  \end{scope}

\end{tikzpicture}
\end{center}


\subsection*{1.6\quad Stabkräfte $S_8$ und $S_9$ -- Rundschnitt am Knoten B}

Knoten B: Stab~[8] (B$\to$E, Richtung $\tfrac{1}{\sqrt{2}}(-1,+1)$),
Stab~[9] (B$\to$F, Richtung $(-1,\,0)$), Auflager $B_V = -\tfrac{P}{2}$.
\begin{align*}
  \sum F_y = 0:\quad & \tfrac{S_8}{\sqrt{2}} + B_V = 0
    &&\Rightarrow\quad S_8 = -\sqrt{2}\,B_V = \tfrac{P}{\sqrt{2}}
       &&\Rightarrow\quad \boxed{S_8 = \tfrac{P}{\sqrt{2}} \approx 14.14\,\text{kN (Zug)}} \\
  \sum F_x = 0:\quad & -\tfrac{S_8}{\sqrt{2}} - S_9 = 0
    &&\Rightarrow\quad S_9 = -\tfrac{S_8}{\sqrt{2}} = -\tfrac{P}{2}
       &&\Rightarrow\quad \boxed{S_9 = -\tfrac{P}{2} = -10.00\,\text{kN (Druck)}}
\end{align*}

\subsection*{1.7\quad Stabkräfte $S_4$, $S_5$, $S_6$ -- Ritterschnitt}

Schnitt durch [4],[5],[6]; rechter Teil \{E,F,B\} ($B_V = -\tfrac{P}{2}$ äußerlich; $P$ liegt im linken Teil).
Stäbe~[4] (C--E) und [6] (D--F): je Neigung $-45^\circ$; Stab~[5] (D--E): horizontal.

\textbf{$S_5$} -- Projektion aller Kräfte des rechten Teils auf $\tfrac{1}{\sqrt{2}}(1,1)$
(senkrecht zu [4] und [6], eliminiert $S_4$ und $S_6$):
\[
  -\tfrac{S_5}{\sqrt{2}} - \tfrac{P}{2\sqrt{2}} = 0
  \quad\Rightarrow\quad \boxed{S_5 = -\tfrac{P}{2} = -10.00\,\text{kN (Druck)}}
\]

\textbf{$S_4$} -- Momentensumme um D
([5] und [6] laufen durch D; nur $S_4$ und $B_V$ liefern Beiträge):
\[
  \sum M_D = 0:\quad \tfrac{S_4}{\sqrt{2}}\cdot\tfrac{L}{2} - \tfrac{P}{2}\cdot L = 0
  \quad\Rightarrow\quad \boxed{S_4 = \sqrt{2}\,P = 28.28\,\text{kN (Zug)}}
\]

\textbf{$S_6$} -- Momentensumme um E
([4] und [5] laufen durch E; nur $S_6$ und $B_V$ liefern Beiträge):
\[
  \sum M_E = 0:\quad -\tfrac{L}{2}\cdot\tfrac{S_6}{\sqrt{2}} - \tfrac{P}{2}\cdot\tfrac{L}{2} = 0
  \quad\Rightarrow\quad \boxed{S_6 = -\tfrac{P}{\sqrt{2}} = -14.14\,\text{kN (Druck)}}
\]

\medskip
\begin{center}
\begin{tabular}{ccccc}
\toprule
$S_8$ & $S_9$ & $S_4$ & $S_5$ & $S_6$ \\
\midrule
$\tfrac{P}{\sqrt{2}} = 14.14$ & $-\tfrac{P}{2} = -10.00$ &
$\sqrt{2}\,P = 28.28$ & $-\tfrac{P}{2} = -10.00$ & $-\tfrac{P}{\sqrt{2}} = -14.14$ \\
\multicolumn{5}{c}{\small [kN];\quad $+$ = Zug,\quad $-$ = Druck} \\
\bottomrule
\end{tabular}
\end{center}

\newpage
\section*{2.~Beispiel \hfill (30\,\%)}

\textbf{Gegeben:} Einfeldträger (Spannweite $L$) mit Dreieckslast
(0 bei A bis $q_0$ bei B),
$A_V = \tfrac{1}{6}q_0 L = 2.00$\,kN,
$B_V = \tfrac{1}{3}q_0 L = 4.00$\,kN, und
\[
  M(x) = -\tfrac{q_0}{6L}\,x^3 + \tfrac{q_0 L}{6}\,x.
\]

\textbf{Systemwerte:} $L = 2$\,m,\quad $q_0 = 6$\,kN/m.

\subsection*{2.1\quad Querkraft $V(x)$}

Durch Ableitung des Momentenverlaufs:
\[
  V(x) = \frac{\mathrm{d}M}{\mathrm{d}x}
        = -\frac{q_0}{2L}\,x^2 + \frac{q_0 L}{6}
  \quad\Rightarrow\quad
  \boxed{V(x) = \frac{q_0 L}{6} - \frac{q_0}{2L}\,x^2}
\]
Kontrolle: $V(0) = \tfrac{1}{6}q_0 L = A_V$\,\checkmark;\quad
$V(L) = \tfrac{1}{6}q_0 L - \tfrac{1}{2}q_0 L = -\tfrac{1}{3}q_0 L = -B_V$\,\checkmark.

\subsection*{2.2\quad Ort des maximalen Moments ($V = 0$)}

\[
  \frac{q_0 L}{6} = \frac{q_0}{2L}\,x^2
  \;\Rightarrow\; x^2 = \frac{L^2}{3}
  \;\Rightarrow\; \boxed{x_{\max} = \frac{L}{\sqrt{3}} \approx 1.1547\,\text{m}}
\]

\subsection*{2.3\quad Maximales Moment}

\[
  M_{\max} = M\!\left(\tfrac{L}{\sqrt{3}}\right)
  = -\frac{q_0}{6L}\cdot\frac{L^3}{3\sqrt{3}} + \frac{q_0 L}{6}\cdot\frac{L}{\sqrt{3}}
  = \frac{q_0 L^2}{6}\cdot\frac{2}{3\sqrt{3}}
  \;\Rightarrow\;
  \boxed{M_{\max} = \frac{q_0 L^2}{9\sqrt{3}}
         = \frac{\sqrt{3}}{27}\,q_0 L^2 \approx 1.5396\,\text{kNm}}
\]

\subsection*{2.4\quad Verläufe}


\begin{center}
\begin{tikzpicture}[xscale=3.5, yscale=3.5, >=stealth, thick, font=\small]

  %====================================================================
  % SHEAR FORCE  V(x)
  %====================================================================
  \begin{scope}[yshift=0cm]
    \draw[->] (-0.05,0) -- (1.1,0) node[right]{$x$};
    \node[left]  at (0,0)  {A};
    \node[right] at (1,0)  {B};

    \filldraw[fill=red!15, draw=red!60!black, thick]
      (0,0) -- plot coordinates { (0.00000,0.27500) (0.02500,0.27448) (0.05000,0.27294) (0.07500,0.27036) (0.10000,0.26675) (0.12500,0.26211) (0.15000,0.25644) (0.17500,0.24973) (0.20000,0.24200) (0.22500,0.23323) (0.25000,0.22344) (0.27500,0.21261) (0.30000,0.20075) (0.32500,0.18786) (0.35000,0.17394) (0.37500,0.15898) (0.40000,0.14300) (0.42500,0.12598) (0.45000,0.10794) (0.47500,0.08886) (0.50000,0.06875) (0.52500,0.04761) (0.55000,0.02544) (0.57500,0.00223) (0.60000,-0.02200) (0.62500,-0.04727) (0.65000,-0.07356) (0.67500,-0.10089) (0.70000,-0.12925) (0.72500,-0.15864) (0.75000,-0.18906) (0.77500,-0.22052) (0.80000,-0.25300) (0.82500,-0.28652) (0.85000,-0.32106) (0.87500,-0.35664) (0.90000,-0.39325) (0.92500,-0.43089) (0.95000,-0.46956) (0.97500,-0.50927) (1.00000,-0.55000) } -- (1,0) -- cycle;

    \node[right, red!70!black] at (0,  0.2750) {$+2.00$\,kN};
    \node[right, red!70!black] at (1,  -0.5500) {$-4.00$\,kN};
    \draw[thin,dashed,gray] (0.5774,0) -- (0.5774,0.0000);
    \node[above, gray, font=\footnotesize] at (0.5774, 0)
          {$x_{max}=\tfrac{L}{\sqrt{3}}$};
    \node[rotate=90, red!70!black] at (-0.22, 0.1) {$V(x)$};
  \end{scope}

  %====================================================================
  % MOMENT  M(x)
  %====================================================================
  \begin{scope}[yshift=-1.9cm]
    \draw[->] (-0.05,0) -- (1.1,0) node[right]{$x$};
    \node[left]  at (0,0) {A};
    \node[right] at (1,0) {B};

    \filldraw[fill=blue!15, draw=blue!60!black, thick]
      (0,0) -- plot coordinates { (0.00000,0.00000) (0.02500,0.03570) (0.05000,0.07127) (0.07500,0.10657) (0.10000,0.14147) (0.12500,0.17583) (0.15000,0.20952) (0.17500,0.24241) (0.20000,0.27436) (0.22500,0.30524) (0.25000,0.33491) (0.27500,0.36324) (0.30000,0.39010) (0.32500,0.41535) (0.35000,0.43886) (0.37500,0.46050) (0.40000,0.48012) (0.42500,0.49761) (0.45000,0.51281) (0.47500,0.52561) (0.50000,0.53585) (0.52500,0.54342) (0.55000,0.54818) (0.57500,0.54999) (0.60000,0.54871) (0.62500,0.54423) (0.65000,0.53639) (0.67500,0.52507) (0.70000,0.51013) (0.72500,0.49144) (0.75000,0.46887) (0.77500,0.44228) (0.80000,0.41154) (0.82500,0.37650) (0.85000,0.33705) (0.87500,0.29304) (0.90000,0.24435) (0.92500,0.19083) (0.95000,0.13236) (0.97500,0.06879) (1.00000,0.00000) } -- (1,0) -- cycle;

    % Max moment marker
    \draw[thin, dashed, gray] (0.5774, 0) -- (0.5774, 0.5500);
    \node[below, blue!70!black] at (0.5774, 0.5500)
          {$M_{max}\approx 1.5396$\,kNm};
    \node[rotate=90, blue!70!black] at (-0.22, 0.2750) {$M(x)$};
  \end{scope}

\end{tikzpicture}
\end{center}

\end{document}
